\section{Fazit}\label{fazit}

Die vorliegende Arbeit hat einen zweistufigen hierarchischen \gb{}-Klassifikator auf dem
CICIoT2023-Datensatz~\cite{neto2023ciciot2023} implementiert, evaluiert und mit einer
\shap{}-basierten Erklaerbarkeitskomponente versehen.

\textbf{Kernergebnisse.}
Die Ablation Study liefert ein klar messbares Ergebnis: \pca{} reduziert die \ba{} von
Stufe~2 um 8{,}21~Prozentpunkte (Pipeline~B: 0{,}6301 gegenueber Pipeline~A: 0{,}5480
mit Standardparametern).
Gleichzeitig trainiert Pipeline~B ohne \pca{} um 583~Sekunden schneller.
Die Grid-Search-Optimierung verbessert die \ba{} von Stufe~2 in Pipeline~B um weitere
1{,}78~Prozentpunkte auf 0{,}6479 (auf einem Teilmengedatensatz von 8.045~Instanzen).
Diese Ergebnisse bestaetigen die Ausgangshypothese: \pca{} beeintraechtigt die
Mehrklassenklassifikation messbar und gleichzeitig die semantische Interpretierbarkeit
der \shap{}-Erklaerungen, da Hauptkomponenten keine direkte Entsprechung zu
Netzwerkmerkmalen haben~\cite{alharby2025ciciot}.

\textbf{Wissenschaftlicher Beitrag.}
Gegenueber Raturi et al.\ \cite{raturi2026exploratory} schliesst die vorliegende Arbeit
die identifizierte Forschungsluecke der fehlenden Erklaerbarkeit durch die Integration von
\shap{} als post-hoc-Erklaerungskomponente.
Gegenueber Alharby~\cite{alharby2025ciciot} ergaenzt die Arbeit die \pca{}-basierte
Klassifikation um eine systematische Ablation Study einschliesslich des Einflusses auf die
\shap{}-Feature-Importance-Ranglisten.
Die Kombination aus zweistufiger hierarchischer Architektur, \ba{} als primaerer Metrik
und \shap{}-Analyse auf dem CICIoT2023-Datensatz stellt nach Kenntnis des Autors eine
eigenstaendige Kombination in der Literatur dar.

\textbf{Ausblick.}
Die vollstaendige \shap{}-Analyse auf dem Gesamtmodell nach Reentrainement mit optimierten
Hyperparametern und die Bewertung der Erklaerungsqualitaet nach Hermosilla et al.\
\cite{hermosilla2025xai_forensic} (Fidelitaet, Konsistenz, Jaccard-Aehnlichkeit) sind
naechste Schritte.
Offene Forschungsfragen bleiben die Latenzmessung fuer den Echtzeiteinsatz in
ressourcenbeschraenkten \iot{}-Umgebungen~\cite{raturi2026exploratory} sowie die
Generalisierbarkeit auf weitere \iot{}-Datensaetze~\cite{hamedani2025iotsurvey}.
